\scanpage{84}\mspage{13}
\handwritten{Malheureusement, on ne peut en conclure $\lambda v=\mu w$, p. ex. si on a $\lambda'v=\mu'w$, on en conclut $\lambda v\sim\mu w$, mais nullement $\lambda v=\mu w$. Plus précisément, dès qu'on se donne dans $X$ des flèches $v,w$ \deleted{\unsure{avec}}}
\handwritten{\diagram{\begin{tikzcd}[ampersand replacement=\&,row sep=small]
 \&a\\ x'\arrow[ur,"v"]\arrow[dr,"w"']\&\\ \&b
\end{tikzcd}}}
\handwritten{avec $\lambda'v=\mu'\!w$, on en conclut dans $\widetilde X$ un $\tilde x$ avec}
\handwritten{\diagram{\begin{tikzcd}[ampersand replacement=\&,row sep=small]
 \&\tilde a\\ \tilde x\arrow[ur,"\tilde v"]\arrow[dr,"\tilde w"']\&\\ \&\tilde b
\end{tikzcd}}}
\handwritten{, mais on n'aura pas pour \unsure{autant} \deleted{\unsure{on}} $\lambda v=\mu w$ : ainsi, dans le diagramme \deleted{\unsure{de la|de}} \deleted{page} 7, on \inserted{savait déjà qu'on} peut trouver $v,w$ de façon à rendre \uline{un} des carrés pertinents commu (disons $\lambda'v=\mu'w$), \unsure{ce|du|de} qui assure déjà $\widetilde X_{//(\tilde a,\tilde b)}\ne\varnothing$ --}
\handwritten{mais je ne vois pas de moyen d'assurer en $=$ temps que l'autre le soit.}
% letterhead
MINISTERE DE L'EDUCATION NATIONALE\\
ACADEMIE D'AIX-MARSEILLE\\
\rule{2cm}{0.4pt}\\
Inspection Académique\\
de Vaucluse\\
\rule{2cm}{0.4pt}\\
Tél. : 90 82 99 76\\
Référence : Secrétariat n\textsuperscript{o}

AVIGNON, le\\[1ex]
L'Inspecteur d'académie,\\
directeur des services départementaux\\
de l'Education Nationale de Vaucluse,\\[1ex]
à \ M
% archive sheet number: 83

